% === Begin Label Data ===
% ID: 214
% 难度星级: 
% 标签: 
% 备注: 
% 组卷引用次数: 0
% === End  Label Data ===

\begin{problem}{2026}{G}{上海卷（理）}{23}{数列}
若无穷数列 $\{a_n\}$ 满足：只要 $a_p = a_q$（$p,q \in \mathbb{N}^*$），必有 $a_{p+1} = a_{q+1}$，则称 $\{a_n\}$ 具有性质 $P$。

（1）若 $\{a_n\}$ 具有性质 $P$，且 $a_1 = 1, a_2 = 2, a_4 = 3, a_5 = 2$，$a_6 + a_7 + a_8 = 21$，求 $a_3$；

（2）若无穷数列 $\{b_n\}$ 是等差数列，无穷数列 $\{c_n\}$ 是公比为正数的等比数列，$b_1 = c_5 = 1$，$b_5 = c_1 = 81$，$a_n = b_n + c_n$，判断 $\{a_n\}$ 是否具有性质 $P$，并说明理由；

（3）设 $\{b_n\}$ 是无穷数列，已知 $a_{n+1} = b_n + \sin a_n$（$n \in \mathbb{N}^*$）．求证：“对任意 $a_1$，$\{a_n\}$ 都具有性质 $P$” 的充要条件为 “$\{b_n\}$ 是常数列”．
\end{problem}